\documentclass[french,a4paper]{article}

\usepackage[utf8]{inputenc}
\usepackage[T1]{fontenc}
\usepackage{lmodern}
\usepackage{geometry}
\usepackage{graphicx}
\usepackage{babel}
\usepackage{amsmath}
\usepackage{amsthm}
\usepackage{amssymb}
\usepackage{hyperref}
\usepackage{stmaryrd}
\usepackage{enumitem}
\usepackage{tcolorbox}
\usepackage{dsfont}
\usepackage{multirow}
\usepackage{etoc}
\usepackage{titlesec}
\usepackage{esint}
\usepackage{cancel}
\usepackage{mathdots}

\setlist[itemize]{label=$\bullet$}


\renewcommand{\P}{\mathbb{P}}
\newcommand{\N}{\mathbb{N}}
\newcommand{\Z}{\mathbb{Z}}
\newcommand{\Q}{\mathbb{Q}}
\newcommand{\R}{\mathbb{R}}
\newcommand{\C}{\mathbb{C}}
\newcommand{\K}{\mathbb{K}}
\renewcommand{\L}{\mathbb{L}}
\newcommand{\U}{\mathbb{U}}
\newcommand{\st}{^*}
\newcommand{\tq}{\middle|}
\newcommand{\inv}{^{-1}}
\newcommand{\E}{\mathbb{E}}
\newcommand{\F}{\mathbb{F}}
\newcommand{\V}{\mathbb{V}}
\renewcommand{\ker}{\text{Ker}}
\newcommand{\im}{\text{Im}}
\newcommand{\card}{\text{Card}}
\newcommand{\Tr}{\text{Tr}}

\theoremstyle{plain}
\newtheorem{thm}{Théorème}
\newtheorem*{ind}{Indication}

\theoremstyle{definition}
\newtheorem{dfn}{Definition}

\theoremstyle{remark}
\newtheorem*{rmq}{Remarque}
\newtheorem*{app}{Application}

\newtcolorbox[auto counter]{ex}[2][]{colback=white, colframe=green!50!white, coltitle=black, fonttitle=\bfseries, title={Exercice~\thetcbcounter~: #2}, after title={\hfill #1}}

\title{TD Suites et séries de fonctions}

\begin{document}

\maketitle

\section{Suites de fonctions, modes de convergence}

\begin{ex}[Moulin.01]{}
	Étudier la convergence simple et la convergence uniforme, sur des ensembles à préciser, des suites de fonctions : 
	\begin{enumerate}
		\item $f_n:x\mapsto\sin(x)e^{-nx}$ ;
		\item $g_n:x\mapsto\displaystyle\frac{x\sqrt{n}}{1+nx^2}$ ;
		\item $h_n:x\mapsto\arctan\displaystyle\left(\frac{n-x}{1+nx}\right)$.
	\end{enumerate}
\end{ex}

\begin{ex}[Moulin.02]{}
	Montrer que la suite de fonctions définie par $\displaystyle f_n(x)=\frac{x^n}{1+x^{2n}}$ converge simplement sur $\R$ vers une limite $f$. Étudier la convergence uniforme de cette suite, puis celle de la série de fonction $f_n-f$.
\end{ex}

\begin{ex}[Moulin.03]{}
	Étude (convergence simple, convergence uniforme sur $\R$, convergence uniforme sur des segments de $\R$) de la suite de fonctions définie par $f_n(x)=\displaystyle\frac{1}{1+(1+nx)^2}$ puis de la série $\sum f_n$.
\end{ex}

\begin{ex}[Moulin.04]{$\Delta$ Une étude de suite de fonctions}
	Convergence simple et convergence uniforme sur tout segment de $\R_+$ de la suite $(f_n)$, où : $$f_n(x)=\begin{cases}
		\displaystyle\left(1-\frac{x}{n}\right)^n&\text{si }0\leqslant x\leqslant n \\ 0&\text{sinon}
	\end{cases}$$ Montrer la convergence uniforme sur $\R_+$.
\end{ex}
\begin{proof}
	Exercice classique. Remarquer que la limite simple est $e^{-x}$. Pour la convergence uniforme sur tout segment $[0,b]$, utiliser la formule de Taylor-Lagrange comme on a obtenu la limite simple avec un équivalent. Enfin, soit $\varepsilon>0$. A partir d'un certain $A\geqslant 0$, on a : $$0\leqslant f_n(x)\leqslant e^{_x}\leqslant \varepsilon$$ Et ensuite, pour $n$ assez grand, OK sur le segment $[0,A]$ par ce qui précède.
\end{proof}

\begin{ex}{Une deuxième étude de suite de fonctions}
	Soit $f\in\mathcal{C}^2(\R,\C)$. Montrer que la suite de fonctions $$f_n:x\mapsto n\left[f\left(x+\frac{1}{n}\right)-f(x)\right]$$ converge uniformément sur tout segment de $\R$.
\end{ex}
\begin{proof}
	On prouve que la limite simple est $f'$. Ensuite, l'inégalité de Taylor-Lagrange permet d'obtenir : $$\lvert f_n(x)-f'(x)\rvert\leqslant\frac{1}{2n}\lVert f''\rVert_{\infty,\left[x,x+\frac{1}{n}\right]}$$ On prend alors $x\in[a,b]$ et alors $f''$ est bornée par $M$ sur $[a,b+1]$ donc $$\forall x\in[a,b],\ \lvert f_n(x)-f'(x)\rvert\leqslant\underbrace{\frac{1}{2n}\lVert f''\rVert_{\infty,\left[a,b+1\right]}}_{\text{indépendant de }x}\xrightarrow[n\to+\infty]{}0$$
\end{proof}
\begin{ex}{Une troisième étude de suite de fonctions}
	Convergence simple et uniforme de $$u_n(x)=\prod_{k=1}^{n}(1+x^k)$$ Où ça ?
\end{ex}
\begin{proof}
	En passant au $\ln$, on obtient un domaine qui vaut $[-1,1[$. En effet, si $x\in]-1,1[$, on a : $$\ln(u_n(x))=\sum_{k=1}^{n}\underbrace{\ln(1+x^k)}_{\sim x^k\text{ sommable}}$$ donc convergence. Ensuite, $$\ln(1-a^n)\leqslant\ln(1+x^n)\leqslant a^n$$ si $x\in[-a,a]$ avec $a\in[0,1[$ donc CVN sur $[-a,a]$.
\end{proof}
\begin{ex}{Une quatrième étude de suite de fonctions}
	Etablir la convergence simple et uniforme sur $\R_+$ de la suite $(f_n)$ de fonctions définies par : $$f_n(t)=\frac{\sin(nt)}{n\sqrt{t}}\ \ \text{si }t>0 \ \ \text{et} \ \ f(0)=0$$
\end{ex}
\begin{proof}
	La limite simple est la fonction nulle. Ensuite, on utilise l'\textbf{arme anti-colleur} (merci pour la dénomination François). On considère $$\varphi(u)=\frac{\sin(u)}{\sqrt{u}}$$ Alors $\varphi$ est bornée sur $\R_+$ (après prolongement par continuité en $0$). Et ensuite : $$\lvert f_n(t)\rvert\leqslant \frac{\mathcal{N}_\infty(\varphi)}{\sqrt{n}}$$ donc on a une convergence normale.
\end{proof}
\begin{rmq}
	\textbf{Quand il y a du $nt$, penser à utiliser cette technique !!!}
\end{rmq}
\begin{ex}{$\Delta$ Composition}
	Soit $\varphi$ une application continue de $[a,b]$ dans $\C$. Montrer que si $(f_n)$ est une suite de fonctions de $\R$ dans $[a,b]$ convergeant uniformément vers $f:\R\to\R$, alors $(\varphi\circ f_n)$ converge uniformément sur $\R$ vers $\varphi\circ f$.
\end{ex}
\begin{proof}
	Revenir à la définition de la convergence uniforme avec les $\varepsilon$ et utiliser l'uniforme continuité de $\varphi$, permise par le théorème de Heine. On prend le $\eta$ qui correspond et on prend $n_0$ à partir duquel $$\lVert f-f_n\rVert_\infty\leqslant\eta$$ avant de conclure avec l'uniforme continuité de $\varphi$.
\end{proof}
\begin{ex}{$\Delta$ Convergence uniforme sur $\R$ d'une suite de polynômes}
	Soit $(P_n)$ une suite de polynômes de $\C[X]$ convergeant uniformément sur $\R$ vers une fonction $f:\R\to\C$. Montrer que $f$ est polynomiale.
\end{ex}
\begin{proof}
	Avec l'inégalité triangulaire, on montre qu'à partir d'un certain rang $n_0$, on a : $$\lVert P_n-P_{n_0}\rVert_\infty\leqslant1$$ Or, les seuls polynômes bornés sont les polynômes constants. On prend donc $c_n\in\C$ tel que $P_n-P_{n_0}=c_n$. En particulier, $c_n=P_n(0)-P_{n_0}(0)$ Mais par convergence simple vers $f$, on a $$c_n\xrightarrow[n\to+\infty]{}f(0)-P_{n_0}(0)$$ Et avec la limite simple on a : $$f=P_{n_0}+f(0)-P_{n_0}(0)$$ donc $f$ est bien polynomiale.
\end{proof}
\begin{ex}{$\star$ Itérées d'une fonction qui diminue strictement la norme}
	Soit $A$ un réel strictement positif et $f:\R\to\R$ une fonction continue vérifiant $$\forall x\in\R\st,\ \lvert f(x)\rvert<\lvert x\rvert$$ Montrer que la suite $f^=f\circ \cdots\circ f$ converge uniformément vers $0$ sur tout segment $[-A,A]$.
\end{ex}
\begin{proof}
	On remarque par passage à la limite que $f(0)=0$ puis on prend $\varepsilon>0$ suffisamment petit. On applique le théorème des bornes atteintes à $\displaystyle\frac{\lvert f(x)\rvert}{\lvert x\rvert}$ sur $[-A,-\varepsilon]$ et su $[\varepsilon,A]$ pour obtenir qu'elle est contractante (ou presque) sur ces ensembles. Nécessairement, les itérées rentrent dans $[-\varepsilon,\varepsilon]$ à partir d'un certain $n$ qui ne dépend que $A$ et $\varepsilon$.
\end{proof}
\begin{ex}{Une intégrale}
	Soit $f:[0,1]^2\to\R$ continue.
	\begin{enumerate}
		\item Montrer que pour tout suite convergente $(x_n)$ d'éléments de $[0,1]$, la suite de fonctions $(y\mapsto f(x_n,y))$ converge uniformément sur $[0,1]$.
		\item Montrer que $x\mapsto\displaystyle\int_{0}^{1}f(x,y)\ dy$ est continue.
	\end{enumerate}
\end{ex}
\begin{proof}
	On utilise le théorème de Heine pour obtenir que $f$ est uniformément continue et on prend $\varepsilon>0$ et $\eta>0$ associé.
	\begin{enumerate}
		\item On note $x$ la limite. Alors on prend $n_0$ tel que pour $n\geqslant n_0$ on ait $\lvert x_n-x\rvert\leqslant\eta$ puis on utilise l'uniforme continuité.
		\item On fait la différence des deux intégrales et on utilise encore l'uniforme continuité.
	\end{enumerate}
\end{proof}
\begin{ex}{$\Delta$ Limite simple de fonctions lipschitziennes de même rapport sur un segment}
	Soit $k\in\R_+\st$ et $(f_n)_{n\in\N}$ une suite de fonctions toutes $k$-lipschitziennes définies sur un segment $[a,b]$ de $\R$ et à valeurs dans $\C$. On suppose que $(f_n)_{n\in\N}$ converge simplement vers une limite $f$. Montrer que la convergence est uniforme.
\end{ex}
\begin{proof}
	Un passage à la limite montre que $f$ est $k$-lipschitzienne. Ensuite, pour $\varepsilon>0$, on prend une subdivision $a_0<\ldots<a_p$ de $[a,b]$ telle que $$\lvert a_{i+1}-a_i\rvert\leqslant\frac{\varepsilon}{3k}$$ Ensuite, on prend un $n_0\in\N$ tel que $$\forall n\geqslant n_0,\ \lvert f_n(a_i)-f(a_i)\rvert\leqslant\frac{\varepsilon}{3}$$ Finalement, pour $x\in[a,b]$, dont le plus proche élément de la subdivision est $a_i$, on écrit : $$\lvert f_n(x)-f(x)\rvert\leqslant\lvert f_n(x)-f_n(a_i)\rvert+\lvert f_n(a_i)-f(a_i)\rvert+\lvert f(a_i)-f(x)\rvert$$ et on utilise la lipschitzianité et le fait que $$\lvert x-a_i\rvert\leqslant \frac{\varepsilon}{3k}$$ puis l'hypothèse sur $n_0$ pour majorer le tout par $\varepsilon$, indépendamment de $x$.
\end{proof}
\begin{rmq}
	Cette technique de prendre une subdivision est fortement utile.
\end{rmq}
\begin{ex}{$\star$ $\star$ $\star$ Version faible du théorème d'Ascoli}
	\begin{enumerate}
		\item Soit $k>0$ et $(f_n)$ une suite de fonctions de $[0,1]$ dans $[-1,1]$ toutes $k$-lipschitziennes. Montrer que l'on peut extraire de $(f_n)$ une sous-suite convergeant uniformément.
		\item Application :montrer qu'est de dimension finie tout sous-espace vectoriel $F$ de $\mathcal{C}^1([0,1],\R)$ sur lequel $f\mapsto\mathcal{N}_\infty(f)$ et $f\mapsto\mathcal{N}_\infty(f)+\mathcal{N}_\infty(f')$ sont équivalentes.
		\item Application 2 : soit $(f_n)$ une suite de fonctions convexes bornées sur un intervalle $I$ ouvert (donc continues) : montrer que si la suite $(f_n)$ est bornée pour $\mathcal{N}_\infty$, elle admet une sous-suite convergeant uniformément sur tout segment de $I$.
	\end{enumerate}
\end{ex}
\begin{proof}
	Exercice \textbf{très} difficile !
	\begin{enumerate}
		\item Mettre en place une extraction diagonale après avoir pris des limites simples en tout point de $[0,1]\cap \Q$ qui est dénombrable, puis utiliser la lipschitzianité un peu comme dans l'exercice précédent. Pour plus de détail, se référer au DM qui concerne ce théorème.
		\item Utilise l'équivalence de ces normées pour essayer de se ramener à des fonctions lipschitziennes
		\item Écrire une inégalité des pentes et ce que signifie le fait d'être bornée.
	\end{enumerate}
\end{proof}
\begin{ex}{$\star$ Lemme de Riemann-Lebesgue sur un intervalle quelconque}
	On rappelle que si $f$ est continue par morceaux sur un segment $I$, alors $\displaystyle\in_If(t)e^{int}\ dt$ tend vers $0$ lorsque $n$ tend vers $+\infty$. Généraliser ce résultat à un intervalle $I$ quelconque pour une fonction $f$ intégrable sur $I$. On pourra se contenter de traiter le cas $I=\R_+$.
\end{ex}
\begin{proof}
	On définit $$I_n=\int_{0}^{+\infty}f(t)e^{int}\ dt \ \ \text{et} \ \ I_n(p)=\int_{0}^{p}f(t)e^{int}\ dt$$ Alors, on a : $$\lvert I_n-I_n(p)\rvert\leqslant\int_{p}^{+\infty}\lvert f(t)\rvert \lvert e^{int}\rvert\ dt=\underbrace{\int_{p}^{+\infty}\lvert f(t)\rvert \ dt}_{\text{indépendant de }n}\xrightarrow[p\to+\infty]{}0$$ On a donc une convergence uniforme puis on applique le théorème de la double limite en utilisant le cas des segments.
\end{proof}
\section{Théorèmes de Dini}
\begin{ex}{$\Delta\star$ Premier théorème de Dini}
	On considère une suite décroissante de fonctions $(f_n)$ sur un compact $K$ qui converge simplement vers $0$. Alors la convergence est uniforme.
\end{ex}
\begin{proof}
	Déjà, remarquons que toutes les $f_n$ sont par conséquent positives. Il est possible de démontrer le résultat à la main en prenant les \textit{maxima} sur le segment, mais il y a plus élégant. Soit $\varepsilon>0$. On considère $$K_n:=\left\{x\in K \ : \ f_n(x)\geqslant\varepsilon\right\}$$ puis on raisonne par l'absurde. On suppose que tous ces compacts (fermés dans un compact) sont non vides. Alors, puisque $(K_n)$ forme une suite décroissante de compacts, d'après le théorème des compacts emboîtés, on peut prendre $x$ dans l'intersection de tous ces compacts. Mais alors il ne peut y avoir convergence simple en $x$. Cela est absurde. Ainsi, il existe un $K_{n_0}$ vide, et par décroissance, tous les suivants le sont. On a donc bien convergence uniforme.
\end{proof}
\begin{rmq}
	On peut aussi le faire à la main mais c'est un peu bourbier. Sinon, ceux qui veulent pourront le faire avec Borel-Lebesgue mais c'est un peu overkill (et pas démontrable très vite).
\end{rmq}
\begin{ex}{$\star$ Deuxième théorème de Dini}
	On considère une suite de fonctions $(f_n)$ définies sur $[a,b]$ à valeurs réelles qui sont toutes croissantes et qui converge simplement vers $f$ continue. Montrer que
	la convergence est uniforme.
\end{ex}
\begin{proof}
	Un passage à la limite montre que $f$ est croissante. On fixe $\varepsilon>0$ Et par croissance et uniforme continuité de $f$, on peut fixer $c_0<\ldots<c_n$ une subdivision de $[0,1]$ telle que $$\forall i,\ f(c_{i+1})-f(c_i)\leqslant\frac{\varepsilon}{2}$$ Ensuite, si $x\in[0,1]$, on prend $i$ tel que $c_i\leqslant x\leqslant c_{i+1}$. Par croissance de $f$ et de $f_n$ on a alors : $$\begin{cases}
		f_n(x)-f(x)\leqslant f_n(c_{i+1})-f(c_i)\leqslant\frac{\varepsilon}{2}+f_n(c_{i+1})-f(c_{i+1})\\
		f_n(x)-f(x)\geqslant f_n(c_i)-f(c_{i+1})\geqslant f_n(c_i)-f(c_i)-\frac{\varepsilon}{2}
	\end{cases}$$ Or, il y a un nombre fini de $c_i$ donc on peut prend $n_0$ tel que $$\forall i,\ \forall n\geqslant n_0,\ \lvert f_n(c_i)-f(c_i)\rvert\leqslant\frac{\varepsilon}{2}$$ On a donc : $$\forall x,\forall n\geqslant n_0,\ \lvert f_n(x)-f(x)\rvert\leqslant\varepsilon$$ ce qui est bien la convergence uniforme souhaitée.
\end{proof}
\begin{ex}{$\star$ Théorème de Dini pour les fonctions convexes}
	Montrer que si une suite de fonctions convexes sur un intervalle ouvert $I$ converge simplement sur $I$, alors la convergence est uniforme sur tout segment de $I$.
\end{ex}
\begin{proof}
	On note $(f_n)$ la suite de fonctions et $f$ la limite simple. On prend $[b,c]\in I$ puis $(a,d)\in I^2$ tel que $a<b<c<d$. Soit $x\ne y\in[b,c]$. La convexité fournit : $$\frac{f_n(b)-f_n(a)}{b-a}\leqslant\frac{f_n(y)-f_n(x)}{y-x}\leqslant\frac{f_n(d)-f_n(c)}{d-c}$$ Or, les deux suite extrémales admettent des limites indépendantes de $x$ et $y$. Donc on en déduit que $$\forall n,\ \forall x\ne y,\ \left\lvert \frac{f_n(y)-f_n(x)}{y-x}\right\rvert\leqslant M$$ où $M$ ne dépend pas ni de $x$, ni de $y$, ni de $n$. Ensuite, il s'agit d'un des exercices précédents (à refaire donc).
\end{proof}
\begin{rmq}
	il est peut-être possible de se ramener à une autre théorème de Dini, peut-être celui sur les fonctions croissantes.
\end{rmq}
\section{Propriétés des limites simples et uniformes}
\begin{ex}{Zéros}
	Soit $(f_n)$ une suite de fonctions de $[a,b]$ dans $\R$ convergeant uniformément sur $[a,b]$ vers une fonction continue $f$. Montrer que si chacune des fonctions $f_n$ admet un zéro, il en est de même pour $f$. Donner un contre-exemple si l'on ne suppose qu'une convergence simple.
\end{ex}
\begin{proof}
	Extraire une valeur d'adhérence d'une suite de zéros puis utiliser la convergence uniforme. Comme contre-exemple, considérer : $$fn(x)=x^n-\frac{1}{2}$$ sur $[0,1]$.
\end{proof}
\begin{ex}{Uniforme continuité}
	Montrer que si $(f_n)$ converge uniformément sur $A$ vers $f$, et si pour tout $n$, la fonction $f_n$ est uniformément continue sur $A$, alors $f$ l'est également.
\end{ex}
\begin{proof}
	Revenir à la définition. Écrire $$\lVert f(x)-f(y)\rVert\leqslant\lVert f(x)-f_{n_0}(x)\rVert+\lVert f_{n_0}(x)-f_{n_0}(y)\rVert+\lVert f_{n_0}(y)-f_{n_0}(y)\rVert$$ en utilisant les hypothèses de convergence uniforme avec $$\lVert f-f_{n_0}\rVert_\infty\leqslant\frac{\varepsilon}{3}$$ et avec l'hypothèse d'uniforme continuité sur $f_{n_0}$. 
\end{proof}
\begin{ex}{Intégrale et limite nulle}
	Soit $f$ continue par morceaux sur $\R_+$ telle que $\lim_{+\infty}=0$. Montrer que $$\lim_{n\to+\infty}\int_{0}^{1}f(nx)\ dx=0$$ mais qu'il n'y a convergence uniforme vers $0$ de $f_n:x\mapsto f(nx)$ que si $f$ est nulle.
\end{ex}
\begin{proof}
	Revenir avec des $\varepsilon$ et couper l'intégrale en deux en utilisant l'hypothèse. On montre le deuxième point par contraposée en prenant $x_0$ tel que $f(x_0)\ne 0$ et en écrivant $$x_0=n\frac{x_0}{n}$$
\end{proof}
\section{Séries de fonctions}
\begin{ex}{$\star$ Un résultat autour de la constante d'Euler $\gamma$}
	Montrer $$\lim_{x\to 0^+}\left[\left(\sum_{n=1}^{+\infty}\frac{1}{n^{1+x}}\right)-\frac{1}{x}\right]=\gamma$$ où $\gamma$ désigne la constante d'Euler (qui apparaît dans le développement asymptotique de la série harmonique). On pourra remarquer que $$\frac{1}{x}=\lim_{N\to+\infty}\sum_{n=1}^{N}\int_{n}^{n+1}\frac{dt}{t^{1+x}}$$
\end{ex}
\begin{proof}
	On pose $$u_n(x)=\frac{1}{n^{1+x}}-\int_{n}^{n+1}\frac{dt}{t^{1+x}}$$ Un calcul d'intégrale et un DL donnent $$u_n(x)\xrightarrow[x\to 0^+]{}\frac{1}{n}-\ln(n+1)+\ln(n)$$ Or, on sait que : $$\sum_{n=1}^{+\infty}\left(\frac{1}{n}-\ln(n+1)+\ln(n)\right)=\gamma$$ Il suffit alors de montrer la convergence uniforme de $\sum u_n$ au voisinage de $0$. Faire une comparaison série-intégrale pour obtenir : $$\left\lvert\sum_{k=n+1}^{+\infty}\left(\frac{1}{k^{1+x}}-\int_{k}^{k+1}\frac{dt}{t^{1+x}}\right)\right\rvert\leqslant\int_{n}^{n+1}\frac{dt}{t^{1+x}}=\frac{1}{xn^x}\left(1-e^{-x\ln\left(1+\frac{1}{n}\right)}\right)$$ or, pour $n\in\N\st$ et $x\in[0,1]$, on a : $$-x\ln\left(1+\frac{1}{n}\right)\in[-1,0]$$ Avec l'IAF, $\exp$ est $1$-lipschitzienne sur $[-1,0]$ et alors : $$\left\lvert\sum_{k=n+1}^{+\infty}\left(\frac{1}{k^{1+x}}-\int_{k}^{k+1}\frac{dt}{t^{1+x}}\right)\right\rvert\leqslant\frac{\ln\left(1+\frac{1}{n}\right)}{n^x}\leqslant\underbrace{\ln\left(1+\frac{1}{n}\right)}_{\text{indépendant de }x}\xrightarrow[n\to+\infty]{}0$$ On a donc la CVU souhaitée et le théorème d'interversion de limites fournit le résultat.
\end{proof}
\begin{ex}{Suite décroissante positive}
	Soit $(a_n)$ une suite décroissante positive et $u_n:t\mapsto a_nt^n(1-t)$.
	\begin{enumerate}
		\item Montrer que $\sum u_n$ converge simplement sur $[0,1]$.
		\item Montrer que $\sum u_n$ converge normalement sur $[0,1]$ si, et seulement si, $\sum\displaystyle\frac{a_n}{n}$ converge.
		\item $\star$ Montrer que $\sum u_n$ converge uniformément sur $[0,1]$ si, et seulement si, $(a_n)$ tend vers $0$.
	\end{enumerate}
\end{ex}
\begin{proof}
	La troisième question est difficile. Remarquons déjà que $a_n$ est bornée.
	\begin{enumerate}
		\item Si $t\in[0,1[$, Ok car $(a_n)$ est bornée. Si $t=1$, c'est la série nulle.
		\item On étudie les variations de $u_n$ : $$u_n'(t)=a_nt^{n-1}(n-(n+1)t)$$ Cela permet de montrer que $u_n$ admet pour maximum $$a_n\frac{n^n}{(n+1)^{n+1}}$$ On prouve alors aisément que $$\mathcal{N}_\infty(u_n)\sim\frac{a_n}{ne}$$ Par équivalence et comme $1/e\ne 0$, on a le résultat.
		\item Supposons $a_n\to l>0$. Alors on montre aisément avec un télescopage que $$\left\lvert\sum_{n=N}^{+\infty}a_nt^n(1-t)\right\rvert\geqslant lt^N$$ donc pas de convergence uniforme car alors $\lVert R_{N-1}\rVert_\infty\geqslant l$. Supposons $a_n\to 0$. Alors : $$\left\lvert\sum_{n=N}^{+\infty}a_nt^n(1-t)\right\rvert\leqslant a_N\sum_{n=N}^{+\infty}(t^n-t^{n+1})\leqslant a_Nt^N\leqslant \underbrace{a_N}_{\text{ind. de }t}\to 0$$ (la deuxième inégalité est valable aussi si $t=1$ même si le télescopage ne tient plus).
	\end{enumerate}
\end{proof}
\begin{ex}{Sommes de puissances}
	\begin{enumerate}
		\item Pour tout $n\in\N\st$, on pose : $$S_n=\sum_{k=0}^{n}\left(1-\frac{k}{n}\right)^n$$ Montrer que la suite $(S_n)_{n\in\N\st}$ converge et préciser sa limite.
		\item En déduire un équivalent simple de la suite $$S_n'=\sum_{j=0}^{n}j^n$$
	\end{enumerate}
\end{ex}
\begin{proof}
	On fait comme d'habitude.
	\begin{enumerate}
		\item On pose $u_k(n)$ qui vaut $\displaystyle\left(1-\frac{k}{n}\right)$ si $k<n$ et $0$ sinon. On a alors $$S_n=\sum_{k=0}^{+\infty}u_k(n)$$ Or, pour tout $k\in\N$, $u_k(n)\xrightarrow[n\to+\infty]{}e^{-k}$ Et on a $$\lvert _k(n)\rvert\leqslant \underbrace{e^{-k}}_{\text{ind. de }n}\text{ sommable}$$ donc on a une convergence normale. On a donc convergence et la limite est $$\frac{1}{1-e^{-1}}$$
		\item Avec un changement d'indice, on montre que $$S_n=\frac{1}{n^n}S_n'$$ On en déduit que $$S_n'\sim\frac{n^n}{1-e^{-1}}$$
	\end{enumerate}
\end{proof}
\begin{ex}{Une étude de série de fonctions}
	On pose $\displaystyle f=\sum_{n=0}^{+\infty}f_n$ où $f_n:t\mapsto te^{-n^2t^2}$.
	\begin{enumerate}
		\item Définition et continuité de $f$ ?
		\item Limite en $+\infty$ ?
		\item Limite en $0$ ?
	\end{enumerate}
\end{ex}
\begin{proof}
	\begin{enumerate}
		\item \begin{itemize}
			\item Définition : OK pour $t=0$, et OK si $t\ne 0$ car on a un $O(e^{-nt^2})$ qui est sommable car $t^2>0$.
			\item Continuité : avec une majoration du reste par une intégrale et un changement de variable, on trouve une CVU sur tout $]-\infty,-a]$ et $[a,+\infty[$ pour tout $a>0$ donc une continuité en tout point de $\R\st$. Remarquer que la continuité en $0$ est l'objet de la troisième question.
		\end{itemize}
		\item On a en particulier convergence uniforme au voisinage de $+\infty$. or, $f_n(t)\xrightarrow[t\to+\infty]{}0$ donc le théorème de la double limite donne $f(t)\xrightarrow[t\to+\infty]{}0$.
		\item Une comparaison série-intégrale (et le dessin qui va avec !) donnent pour $t>0$ : $$\left\lvert f(t)-\underbrace{\int_{0}^{+\infty}te^{-x^2t^2}\ dx}_{=\sqrt{\pi}/2} \right\rvert\leqslant t$$ donc $$f(t)\xrightarrow[t\to0^+]{}\frac{\sqrt{\pi}}{2}$$ On constate qu'il n'y a pas de continuité en $0$ et donc pas de convergence uniforme au voisinage de $0$.
	\end{enumerate}
\end{proof}
\begin{ex}{Une deuxième étude de série de fonctions}
	\begin{enumerate}
		\item Déterminer le domaine de convergence réel de la série de fonctions de terme général $f_n:x\mapsto\displaystyle\frac{xe^{-nx}}{\ln(n)}$. on note $\varphi(x)$ sa somme.
		\item Montrer qu $\varphi$ est continue sur $\R_+\st$.
		\item La série est-elle normalement convergente sur son domaine de définition ?
		\item Déterminer la limite puis un équivalent de $\varphi$ en $0^+$.
	\end{enumerate}
\end{ex}
\begin{proof}
	\begin{enumerate}
		\item Si $x<0$, divergence grossière. Si $x=0$,$f_n(x)=0$. Si $x>0$, $f_n(x)=o(n^{-2})$ donc OK.
		\item $$\lvert f_n(x)\rvert\leqslant \frac{be^{-na}}{\ln(n)}\text{ sommable}$$ si $x\in[a,b]$ avec $a<b$. CVN sur tout segment de $\R_+\st$ donc continuité.
		\item On écrit $$f_n(x)=\frac{\psi(nx)}{n\ln(n)}$$ avec $\psi:t\mapsto te^{-t}$ On a alors : $$\lVert f_n\rVert_\infty=\frac{\lVert\psi\rVert_\infty}{n\ln(n)}$$ Or, $\lVert\psi\rVert_\infty\ne0$ et la série des inverses de $n\ln(n)$ diverge (cas limite des séries de Bertrand) donc on n'a pas convergence normale sur $\R_+$.
		\item Comparaison série intégrale puis ???
	\end{enumerate}
\end{proof}
\begin{ex}{Une troisième étude de séries de fonctions : $\zeta$ alternée}
	\begin{enumerate}
		\item Ensemble de définition de $\chi:s\mapsto\displaystyle\sum_{n=1}^{+\infty}\frac{(-1)^{n-1}}{n^s}$
		\item Montrer que $\chi$ est de classe $\mathcal{C}^\infty$ sur son ensemble de définition.
		\item Préciser la monotonie de $\chi$ au voisinage de $+\infty$.
		\item $\star$ Montrer que $\chi$ tend vers $1/2$ en $0^+$.
	\end{enumerate}
\end{ex}
\begin{proof}
	La dernière question est difficile.
	\begin{enumerate}
		\item Si $s\leqslant0$, divergence grossière. Si $s>0$, OK avec le théorème des séries alternées.
		\item Dériver et dominer.
		\item Dériver $\chi$ et utiliser le théorème des séries alternées.
		\item Calculer sur $2f$ en sortant un terme d'une des deux sommes et en faisant un décalage d'indice sur l'autre. On utilise ensuite le théorème des séries alternées sur cette série alternée convexe.
	\end{enumerate}
\end{proof}
\begin{ex}{$\star$ Théorème de Mertens}
	Soit $\sum u_n$ et $\sum v_n$ deux séries à termes complexes convergentes. On pose $$w_n=\sum_{p+q=n}u_p v_q$$ Montrer que si $\sum u_n$ est absolument convergente, alors le produit de Cauchy $\sum w_n$ converge et vérifie : $$\sum_{n=0}^{+\infty}w_n=\left(\sum_{n=0}^{+\infty}u_n\right)\left(\sum_{n=0}^{+\infty}v_n\right)$$ On pourra écrire $$\sum_{k=0}^{n}w_k=\sum_{p=0}^{+\infty}u_pV_{n-p}$$ où $(V_n)$ est la suite des sommes partielles de $\sum v_n$ convenablement prolongée à $\Z$.
\end{ex}
\begin{proof}
	Comme suggéré, on définit $V_n=0$ pour $n<0$. On veut alors utiliser le théorème de la double limite. On montre alors une convergence normale. On sait que $(V_n)$ est bornée donc $\lvert u_pV_{n-p}\rvert\leqslant C\lvert u_p\rvert$ qui est sommable et indépendante de $n$. On peut alors utiliser le théorème de la double limite sur les sommes partielles.
\end{proof}
\begin{ex}{$\star$ Réciproque d'un résultat classique}
	Soit $D=\{d_n| n\in\N\}$ un ensemble au plus dénombrable. Alors il existe une fonction $f$ croissante dont l'ensemble des points de discontinuité est exactement $D$.
\end{ex} \begin{proof}
	Pour tout $n\in\N$, on pose $$\begin{array}{lrcl}
		f_n:&\R&\longrightarrow&\R \\
		&x&\longmapsto&\begin{cases}
			1/2^n\ \text{si} \ x\geqslant d_n \\
			0\ \text{si} \ x<d_n
		\end{cases}
	\end{array}$$ Alors $f=\sum f_n$ converge simplement. En effet, pour $x\in\R$, on pose $$I_x=\{n\in\N\ : \ d_n\leqslant x\}$$ Alors $f(x)=\sum_{i\in I_x}2^{-i}$ qui est bien finie comme sous-famille de la famille $(2^{-n})_{n\in\N}$ qui est sommable. De plus, $f$ est croissante comme limite simple de fonctions croissantes. Par conséquent, elle admet des limites à gauche et à droite en tout point. Soit $n\in\N$ et notons $l_n$ cette limite. Il est clair que la limite à droite de $f$ en $d_n$ est supérieure à $f(d_n)$ par croissance. Or, pour $x<d_n$ : $$f(x)=\sum_{i\in I_x}2^{-i}\leqslant\sum_{d_i<d_n}2^{-i}=f(d_n)-\frac{1}{2^n}$$ Donc $l_n<f(d_n)$ et $d_n$ est un point de discontinuité de $f$. Et, si $x\notin D$, on peut distinguer les cas : 
	\begin{itemize}
		\item si $x$ n'est pas un point d'accumulation de $D$, alors on a convergence normale au voisinage de $x$ donc continuité car les $f_i$ sont continues au voisinage de $x$ ;
		\item si $x$ est un point d'accumulation de $D$ : on revient à la définition en prenant $\varepsilon$, puis on s'approche suffisamment de $x$ pour que la différence causée par une $\sum_{j\in J}2^{-j}$ soit inférieure à $\varepsilon$ (possible car la série converge). 
	\end{itemize}
\end{proof}
\section{Approximation uniforme}
\begin{ex}{Approximation de la fonction et de ses dérivées sur un segment}
	Soit $f:[0,1]\to\R$.
	\begin{enumerate}
		\item Montrer que si $f$ est de classe $\mathcal{C}^p$ avec $p\in\N$, il existe une suite de polynômes $(P_n)$ telle que $(P_n^{(k)})$ converge uniformément vers $f^{(k)}$ pour tout $k\in\llbracket0,p\rrbracket$.
		\item $\star$ Montrer que si $f$ est de classe $\mathcal{C}^\infty$, alors il existe une suite de polynômes $(P_n)$ telle que $(P_n^{(k)})$ converge uniformément vers $f^{(k)}$ pour tout $k\in\N$.
	\end{enumerate}
\end{ex}
\begin{proof}
	La deuxième question est difficile.
	\begin{enumerate}
		\item Prendre l'exemple de la fin du cours de fonctions vectorielles et ajouter $$P_n=\sum_{k=0}^{p-1}\frac{f^{(k)}(0)}{k!}X^k$$
		\item Soit $p\in\N$. Par la première question, prenons $P_p\in\R[X]$ tel que $$\forall k\leqslant p,\ \lVert P_p^{(k)}-f^{(k)}\rVert_\infty\leqslant\frac{1}{2^p}$$ On montre que $(P_p)_{p\in\N}$ convient. Soit $k\in\N$. Si $p\geqslant k$, on a par construction $$\lVert P_p^{(k)}-f^{(k)}\rVert_\infty\leqslant\frac{1}{2^p}\xrightarrow[p\to+\infty]{}0$$
	\end{enumerate}
\end{proof}
\begin{ex}{$\Delta$ Application du théorème de Weierstrass}
	Soit $f:[a,b]\to\C$ continue telle que $$\forall n\in\N,\ \int_{a}^{b}t^nf(t)\ dt=0$$ Montrer que $f$ est nulle. On pourra commencer par le cas où $f$ est à valeurs réelles.
\end{ex}
\begin{proof}
	Le prouver dans le cas où $f$ est réelle. Le résultat s'en déduire en passant aux parties réelle et imaginaire. utiliser le théorème de Weierstrass pour approximer $f$ par une suite de polynômes. Alors, l'hypothèse permet de montrer que l'intégrale de $f^2$ est nulle, donc $f=0$ car continue positive.
\end{proof}
\begin{ex}{Analogue au précédent}
	Montrer qu'il n'existe pas de fonction $f\in\mathcal{C}^0([0,1])$ telle que $$\int_{0}^{1}xf(x)\ dx=1 \ \ \text{et}\ \ \forall n\geqslant 2,\ \int_{a}^{b}x^nf(x)\ dx=0$$
\end{ex}
\begin{proof}
	Raisonnons par l'absurde et supposons l'existence d'une telle fonction $f$. Alors, en faisant comme dans l'exercice précédent, on montre que $$\forall x\in[0,1],\ x^2f(x)=0$$ On en déduit que $$\forall x\in[0,1],\ xf(x)=0$$ et donc que $$\int_{0}^{1}xf(x)\ dx=0$$ ce qui est absurde.
\end{proof}
\section{TD Châteaux}
\begin{ex}{Théorèmes de Dini}
	\begin{enumerate}
		\item Montrer le premier théorème de Dini : soit $K$ un compact et $(f_n)$ une suite croissante de fonctions continues de $K$ dans $\R$ qui converge simplement vers $f$ continue. Alors, la convergence est uniforme. Indication : on pourra considérer $\Omega_n=\{x\in K\ \mid\ f(x)-\varepsilon< f_n(x)\}$ et utiliser le théorème des compacts emboités.
		\item Montrer le deuxième théorème de Dini : soit $(f_n)$ une suite de fonctions croissantes et continues sur $[a,b]$ qui converge simplement vers $f$ continue. Alors, la convergence est uniforme. Indication : on pourra utiliser le théorème de Heine.
	\end{enumerate}
\end{ex}
\begin{ex}{Convergence d'une suite de fonctions lipschitziennes}
	\begin{enumerate}
		\item Soit $(f_n)$ une suite de fonctions toutes $M$-lipschitziennes sur $[a,b]$ et qui converge simplement vers $f$. Montrer que $f$ est $M$-lipschitzienne et que la convergence est uniforme.
		\item On pose $f_n(x)=\min(nx,\sqrt{x})$. Vérifier que chaque fonction $f_n$ est lipschitzienne. Démontrer que la suite $(f_n)$ converge uniformément sur $[0,1]$ vers la fonction racine carré, qui n'est pas lipschitzienne sur $[0,1]$.
		\item Soit $(f_n)$ une suite de fonctions convexes qui converge simplement vers $f$ sur un intervalle $I$. Démontrer que la convergence est uniforme sur tout segment inclus dans l'intérieur de $I$. 
	\end{enumerate}
\end{ex}
\begin{ex}{Fonctions $\alpha$ et $\zeta$ de Riemann}
	On pose, sous réserve d'existence, $\zeta(s)=\displaystyle\sum_{n=1}^{+\infty}\frac{1}{n^s}$ et $\alpha(s)=\displaystyle\sum_{n=1}^{+\infty}\frac{(-1)^n}{n^s}$.
	\begin{enumerate}
		\item Démontrer la convergence normale sur tout compact du demi-plan $\text{Re}(s)>1$ de ces deux séries de fonctions. En déduire que les fonctions $\zeta$ et $\alpha$ sont définies dans le champ complexe pour $\text{Re}(s)>1$, et sont continues sur ce demi-plan. \\
		Désormais, \textbf{on se limite à la variable réelle}.
		\item Démontrer que $\zeta$ est de classe $\mathcal{C}^\infty$ sur $]1,+\infty[$ et que $\zeta'(x)=\displaystyle-\sum_{n=1}^{+\infty}\frac{\ln(n)}{n^x}$.
		\item Démontrer que la série de fonctions définissant $\alpha$ converge uniformément sur tout segment inclus dans $]0,+\infty[$. En déduire que $\alpha$ est continue sur $]0,+\infty[$, et de classe $\mathcal{C}^\infty$ sur cet intervalle.
		\item Démontrer : $\forall x>1,\ \alpha(x)=(1-2^{1-x})\zeta(x)$. Indication : on pourra séparer les indices pairs et impairs.
		\item Déterminer $\alpha(1)$. En déduire que $\zeta(x)\displaystyle\underset{x\to 1^+}{\sim}\frac{1}{x-1}$.
		\item Démontrer : $\zeta(x)-1\underset{x\to1^+}{\sim}\displaystyle\frac{1}{2^x}$. Indication : on pourra effectuer une comparaison série-intégrale.
	\end{enumerate}
\end{ex}
\begin{ex}{Équicontinuité}
	Soit $K$ une partie compacte d'un espace vectoriel normée $E$. L'espace $\mathcal{E}=\mathcal{C}^0(K,\C)$ est muni de la norme infinie. On dit qu'une partie $\mathcal{F}$ de $\mathcal{E}$ est équicontinue lorsque $$\forall\varepsilon>0,\ \forall x\in K,\ \exists\eta>0,\ \forall y\in K,\ \forall f\in\mathcal{F},\ \lVert x-y\rVert\leqslant\eta\implies\lvert f(x)-f(y)\rvert\leqslant\varepsilon$$ On dit qu'une partie $\mathcal{F}$ de $\mathcal{E}$ est uniformément équicontinue lorsque $$\forall\varepsilon>0,\ \exists\eta>0,\ \forall(x,y)\in K^2,\ \forall f\in\mathcal{F},\ \lVert x-y\rVert\implies\lvert f(x)-f(y)\rvert\leqslant\varepsilon$$
	\begin{enumerate}
		\item Montrer que toute partie équicontinue est uniformément équicontinue.
		\item Soit $(f_n)$ une suite d'éléments de $\mathcal{E}$ qui converge simplement et telle que l'ensemble $\{f_n\ \mid\ n\in\N\}$ soit équicontinu. Montrer que $(f_n)$ converge uniformément vers $f$.
		\item Soit $A$ une partie dense de $K$ et soit $\{f_n\ \mid\ n\in\N\}$ une partie équicontinue de $\mathcal{E}$ telles que $(f_n)$ converge simplement vers $f$ sur $A$. Montrer que $f$ se prolonge en une application $\hat{f}$ continue sur $K$ telle que $(f_n)$ converge simplement vers $\hat{f}$ sur $K$.
		\item \textbf{Théorème d'Ascoli} : supposons $E$ de dimension finie et soit $\mathcal{F}$ une partie de $\mathcal{E}$. Prouver que les assertions suivantes sont équivalentes :
		\begin{enumerate}[label=(\roman*)]
			\item $\mathcal{F}$ est d'adhérence compacte
			\item $\mathcal{F}$ est équicontinue et $\forall x\in K,\ \{f(x)\ \mid \ f\in\mathcal{F}\}$ est borné
		\end{enumerate}
		Indication : pour le sens $(ii)\implies(i)$, on montrera que de toute suite de $\mathcal{F}$, on peut extraire une suite convergeant simplement dans $K\cap D$ où $D$ est dénombrable dense par le biais d'un procédé diagonal.
		\item Application : retrouver le résultat de la première question sur la convergence d'une suite de fonctions toutes lipschitziennes de même rapport.
	\end{enumerate}
\end{ex}
\end{document}